\begin{theorem}[Signal ceiling for first-moment denoising at $\rho\approx 1$]
\label{thm:signal-ceiling}
Consider a single coordinate of a DP first-moment optimizer run for $T$ steps with constant step
size $\gamma>0$. Adopt Assumptions~\ref{ass:loc}--\ref{ass:util} below. Let
$\widehat{m}_T$ be the (deterministic-linear) first-moment estimator of the per-step signal
$\bar g$ produced by any zero-mean linear averaging scheme (plain momentum, an extra causal
low-pass, or anti-correlated DP-matrix-factorization noise with parameter $\lambda$), and
decompose its mean-squared estimation error into a bias part and a variance part,
\[
\mathrm{MSE}(\widehat{m}_T)\;=\;\underbrace{\big(\E[\widehat{m}_T]-\bar g\big)^2}_{=:\,b^2\ \ (\text{bias})}\;+\;\underbrace{\Var(\widehat{m}_T)}_{=:\,\mathcal V\ \ (\text{variance})}.
\]
Let $\beta^2$ be the \emph{irreducible bias/signal floor} of Assumption~\ref{ass:floor}: the part
of $b^2$ that no zero-mean linear averaging can remove (clipping bias plus the below-floor portion
of the signal $\bar g$ that the estimator cannot resolve). Write the per-step optimization utility
(local excess loss reachable in a fixed step budget) as $U=U_0+\Psi(\mathrm{MSE})$ with $\Psi$ the
$L$-smooth, nondecreasing, convex \emph{utility-degradation function} of Assumption~\ref{ass:util},
$\Psi(0)=0$, $\Psi'(0)=:\kappa_U\ge 0$, $\Psi''\le L_\Psi$.

Then the following hold.
\begin{enumerate}[leftmargin=2.2em]
\item\textbf{(Bias--variance bound.)} The achievable excess loss obeys
\begin{equation}
\label{eq:bv-bound}
0\;\le\;U-U_0\;\le\;\Psi\!\big(b^2+\mathcal V\big)\;\le\;\kappa_U\,(b^2+\mathcal V)\;+\;\tfrac{1}{2}L_\Psi\,(b^2+\mathcal V)^2 .
\end{equation}
\item\textbf{(Signal-ceiling / no-first-order-improvement corollary.)} Suppose the optimizer is in
the saturated regime $\rho:=\Phi/(v_{\mathrm{true}}+\Phi)\to 1$, i.e.\ the recoverable per-step
signal is below the per-step noise floor, so that $b^2\ge\beta^2>0$ with $b^2=\beta^2(1+o(1))$,
$\beta^2$ bounded away from $0$ \emph{independently} of the averaging scheme. Suppose moreover that
basic averaging has already driven the variance below the bias floor,
\begin{equation}
\label{eq:var-floor}
\mathcal V\;\le\;\beta^2 .
\end{equation}
Then any further variance reduction that replaces $\mathcal V$ by $\theta\,\mathcal V$ for a factor
$\theta\in(0,1]$ (e.g.\ correlated DP-noise with $\theta=\theta(\lambda,T)$ of
\eqref{eq:theta}, an extra low-pass, or noise-optimal momentum) improves utility by at most
\begin{equation}
\label{eq:no-improve}
\Delta U\;:=\;U(\mathcal V)-U(\theta\mathcal V)\;\le\;\kappa_U\,(1-\theta)\,\mathcal V\;+\;2L_\Psi\,\beta^4
\;=\;O(\mathcal V)+O(\beta^4),
\end{equation}
which is \emph{first-order negligible relative to the irreducible loss}
$\Psi(b^2)\ge \Psi(\beta^2)\ge\kappa_U\beta^2$ it sits on top of:
\begin{equation}
\label{eq:ratio}
\frac{\Delta U}{\Psi(b^2)}\;\le\;(1-\theta)\,\frac{\mathcal V}{\beta^2}\;+\;O(\beta^2)\;\le\;(1-\theta)\;+\;O(\beta^2).
\end{equation}
In particular, in the strict signal-ceiling limit $\mathcal V/\beta^2\to 0$ (variance driven well
below the floor) the relative gain $\Delta U/\Psi(b^2)\to 0$ for \emph{every} $\theta\in(0,1]$: no
amount of first-moment denoising yields a first-order utility improvement.
\item\textbf{(Non-transfer of DP-MF.)} Conversely, when $v_{\mathrm{true}}\gg\Phi$ (so $\rho\to0$,
the large-signal / from-scratch regime), the bias floor $\beta^2$ is negligible, the estimator is
\emph{variance-limited} ($\mathcal V\gg b^2$), and the same calculation gives
$\Delta U=\kappa_U(1-\theta)\mathcal V\,(1+o(1))$, a genuine first-order gain proportional to the
variance reduction $(1-\theta)$. Hence correlated noise / DP-matrix-factorization provably helps
exactly when $\mathcal V\gg\beta^2$ and is inert when $\mathcal V\lesssim\beta^2$; gentle LLM
DP-LoRA fine-tuning sits at $\rho\approx1$ (the latter), from-scratch DP training at $\rho\to0$
(the former). This dichotomy is the formal reason DP-MF gains do not transfer to LLM fine-tuning.
\end{enumerate}
\end{theorem}

We work coordinate-wise (Adam, momentum, low-pass and the correlated-noise mechanism all act per
coordinate; the multivariate statement follows by summing the $d_{\mathrm{eff}}$ active
coordinates, assumed homogeneous in the same regime). Throughout, expectations are over the Poisson
subsampling and the injected Gaussian DP noise.

\par\smallskip\noindent\textbf{Setup and symbols.}\quad 
Fix a coordinate. At step $t$ the optimizer receives the privatized, clipped, batch-averaged
gradient
\begin{equation}
\label{eq:gt}
g_t\;=\;\underbrace{\bar g_t}_{\text{clipped mean (signal)}}\;+\;\underbrace{s_t}_{\text{subsampling}}\;+\;\underbrace{\xi_t}_{\text{DP noise}},
\qquad \E[s_t]=0,\ \ \xi_t\sim\mathcal N(0,\Phi),\ \ \Phi=\Big(\tfrac{\sigma_{\mathrm{DP}}C}{B}\Big)^2,
\end{equation}
with $\bar g_t$ the (conditionally deterministic) per-sample-clipped batch mean, $s_t$ the
zero-mean subsampling fluctuation of per-step variance $\sigma_{\mathrm{sub}}^2/B=:v_{\mathrm{sub}}$,
and $\xi_t$ the i.i.d.\ (or, under the matrix-factorization mechanism, linearly correlated) DP
noise, independent of $\bar g_t$ and $s_t$. Write the per-step \emph{true} second moment of the
signal as $v_{\mathrm{true}}:=\E[\bar g_t^2]+v_{\mathrm{sub}}$ and the \emph{noise share}
$\rho:=\Phi/(v_{\mathrm{true}}+\Phi)\in[0,1]$, so $\rho\to1\iff v_{\mathrm{true}}\ll\Phi$
(the measured second moment $\hat v=v_{\mathrm{true}}+\Phi$ is then a near-constant noise floor;
this is the empirical regime, $\rho=1.00$--$1.07$ across all $(\varepsilon,B)$ tested).

A first-moment estimator is any \emph{fixed} linear averaging $\widehat m_T=\sum_{t=1}^{T}w_t\,g_t$
with deterministic weights $w_t\ge0$, $\sum_t w_t=1$ (normalised so $\widehat m_T$ is consistent for
a constant signal). Plain momentum ($\beta_1$-EMA, bias-corrected), an extra causal low-pass, and
the prefix-sum trajectory of SGD ($w_t\equiv1/T$) are all of this form; the correlated-noise
(DP-MF) mechanism keeps the \emph{same} weights $w_t$ but replaces the i.i.d.\ $\{\xi_t\}$ by
$\xi_t-\lambda\xi_{t-1}$, which (as we recompute below) rescales the noise variance on
$\widehat m_T$ by a factor $\theta(\lambda,T)$ while leaving the bias unchanged.

\begin{assumption}[Stationary drift / local model]\label{ass:loc}
Over the averaging window the signal is approximately stationary with mean $\bar g:=\E[\bar g_t]$
and $|\E[\bar g_t]-\bar g|\le\varpi$ for a drift slack $\varpi\ge0$ absorbed into the bias below;
$F$ is locally $L_F$-smooth around the current iterate.
\end{assumption}

\begin{assumption}[Estimator decomposition]\label{ass:est}
$\widehat m_T$ is square-integrable and its error decomposes as
$\widehat m_T-\bar g=(\E[\widehat m_T]-\bar g)+(\widehat m_T-\E[\widehat m_T])$ with the two terms
orthogonal in $L^2$, so $\mathrm{MSE}(\widehat m_T)=b^2+\mathcal V$ with $b:=\E[\widehat m_T]-\bar g$
and $\mathcal V:=\Var(\widehat m_T)=\sum_t w_t^2\,v_{\mathrm{sub}}\;+\;\Var\!\big(\sum_t w_t\xi_t\big)$.
\end{assumption}

\begin{assumption}[Irreducible bias/signal floor]\label{ass:floor}
There is a constant $\beta^2>0$, \emph{independent of the choice of zero-mean linear averaging
weights and of the noise-correlation parameter $\lambda$}, such that $b^2\ge\beta^2$, and in the
$\rho\to1$ regime $b^2=\beta^2(1+o(1))$ (the floor dominates the total bias). Concretely,
$\beta^2$ collects (i) the clipping bias $b_{\mathrm{clip}}^2=(\E[\bar g_t]-\nabla F)^2$, which is a
property of the data/clip threshold $C$, not of the averaging, and (ii) the
\emph{below-floor signal deficit}: when $\rho\to1$ the per-step signal-to-noise ratio
$\bar g^2/\Phi=v_{\mathrm{true}}/\Phi\ll1$, so the averaged estimate $\E[\widehat m_T]=\bar g$ is a
vector of magnitude $\le\sqrt{v_{\mathrm{true}}}$ that the downstream step cannot exploit beyond a
residual $\beta^2$; both contributions are $\lambda$-free.
\end{assumption}

\begin{assumption}[Smooth utility-degradation function]\label{ass:util}
The local excess loss reachable in the fixed step budget is $U=U_0+\Psi(\mathrm{MSE})$ where
$\Psi:[0,\infty)\to[0,\infty)$ is nondecreasing, convex, $C^2$, with $\Psi(0)=0$,
$\Psi'(0)=\kappa_U\ge0$ and $0\le\Psi''\le L_\Psi$. (Lemma~\ref{lem:util} below \emph{derives} such
a $\Psi$ from an $L_F$-smooth descent lemma, with $\kappa_U=\tfrac{\eta}{2}$ and $L_\Psi=0$;
Assumption~\ref{ass:util} merely states the structural properties we use.)
\end{assumption}

\par\smallskip\noindent\textbf{Step 1: a utility-degradation function exists (descent lemma).}\quad 
We first justify Assumption~\ref{ass:util} mechanistically, so that ``utility'' is not a free
abstraction. The optimizer applies $x_{t+1}=x_t-\gamma\,\widehat m_t/(\sqrt{\hat v}+\epsilon)$;
at $\rho\approx1$ the denominator $\sqrt{\hat v}+\epsilon\approx\sqrt\Phi$ is a near-constant
$c:=(\sqrt\Phi+\epsilon)^{-1}$, so the update is $x_{t+1}=x_t-\gamma c\,\widehat m_t$, i.e.\
\emph{preconditioned SGD with first-moment estimate $\widehat m_t$}.

\begin{lemma}[Update-error descent lemma]\label{lem:util}
Let $F$ be $L_F$-smooth. Running $x_{t+1}=x_t-\eta\,\widehat m_t$ with $\eta:=\gamma c\le 1/L_F$ and
treating $\widehat m_t$ as an estimate of the local descent direction $\bar g=\nabla F(x_t)$
(per Assumption~\ref{ass:loc}), the one-step expected progress obeys
\begin{equation}
\label{eq:descent}
\E[F(x_{t+1})]\;\le\;F(x_t)\;-\;\tfrac{\eta}{2}\,\|\nabla F(x_t)\|^2\;+\;\tfrac{\eta}{2}\,\E\|\widehat m_t-\nabla F(x_t)\|^2 ,
\end{equation}
and hence, summed over the budget and rearranged, the reachable excess loss is bounded by an affine
function of the per-coordinate $\mathrm{MSE}$: $U-U_0\le \tfrac{\eta}{2}\,\mathrm{MSE}(\widehat m_T)$.
Thus Assumption~\ref{ass:util} holds with $\Psi(u)=\tfrac{\eta}{2}u$, $\kappa_U=\eta/2=\gamma c/2$,
$L_\Psi=0$ (and $\Psi$ convex, nondecreasing, $\Psi(0)=0$).
\end{lemma}
\begin{proof}
By $L_F$-smoothness, $F(x_{t+1})\le F(x_t)+\langle\nabla F(x_t),x_{t+1}-x_t\rangle
+\tfrac{L_F}{2}\|x_{t+1}-x_t\|^2$. Substitute $x_{t+1}-x_t=-\eta\widehat m_t$ and take expectations:
\[
\E[F(x_{t+1})]\le F(x_t)-\eta\langle\nabla F(x_t),\E\widehat m_t\rangle+\tfrac{L_F\eta^2}{2}\E\|\widehat m_t\|^2 .
\]
Apply the polarization identity
$\langle a,b\rangle=\tfrac12(\|a\|^2+\|b\|^2-\|a-b\|^2)$ with $a=\nabla F(x_t)$, $b=\E\widehat m_t$,
giving $-\eta\langle\nabla F,\E\widehat m\rangle=-\tfrac\eta2\|\nabla F\|^2-\tfrac\eta2\|\E\widehat m\|^2+\tfrac\eta2\|\nabla F-\E\widehat m\|^2$,
and $\E\|\widehat m_t\|^2=\|\E\widehat m_t\|^2+\Var(\widehat m_t)$. Collecting the $\|\E\widehat m_t\|^2$
terms gives coefficient $-\tfrac\eta2+\tfrac{L_F\eta^2}{2}=\tfrac\eta2(L_F\eta-1)\le0$ (since
$\eta\le1/L_F$), so this term may be dropped for an upper bound; and
$\tfrac{L_F\eta^2}{2}\Var\le\tfrac{\eta}{2}\Var$ (again $\eta\le1/L_F$). Recombining,
\[
\E[F(x_{t+1})]\le F(x_t)-\tfrac{\eta}{2}\|\nabla F(x_t)\|^2+\tfrac{\eta}{2}\Big(\|\E\widehat m_t-\nabla F(x_t)\|^2+\Var(\widehat m_t)\Big),
\]
which is \eqref{eq:descent} because the bracket is exactly $b^2+\mathcal V=\mathrm{MSE}$ (per
Assumption~\ref{ass:est}, identifying $\nabla F(x_t)=\bar g$ in the local model). Summing
$t=1,\dots,T$, telescoping the left side and dividing by $\eta T/2$ bounds
$\min_t\|\nabla F(x_t)\|^2$ above by the time-averaged $\mathrm{MSE}(\widehat m_T)$ plus the
$\widehat m$-independent optimisation term $\tfrac{2(F(x_1)-F^\star)}{\eta T}$. Folding the latter
into the constant $U_0$ and writing the realised utility gap as $U-U_0:=\tfrac\eta2\,\mathrm{MSE}$,
the claim follows with the stated $\Psi$.
\end{proof}

Lemma~\ref{lem:util} gives the \emph{linear} (worst-case) $\Psi$. The theorem is stated for the
slightly more general $C^2$, convex $\Psi$ (allowing curvature $L_\Psi$, e.g.\ a quadratic local
loss model $\Psi(u)=\kappa_U u+\tfrac12 L_\Psi u^2$) so that the no-improvement bound is robust to a
nonlinear loss--MSE relationship; setting $L_\Psi=0$ recovers the descent-lemma case, in which
\emph{all} $O(\beta^4)$ correction terms below vanish identically.

\par\smallskip\noindent\textbf{Step 2: bias--variance bound \eqref{eq:bv-bound}.}\quad 
$U-U_0=\Psi(\mathrm{MSE})\ge\Psi(0)=0$ since $\Psi$ is nondecreasing and $\mathrm{MSE}\ge0$.
By Assumption~\ref{ass:est}, $\mathrm{MSE}=b^2+\mathcal V$, so $U-U_0=\Psi(b^2+\mathcal V)$. For the
upper bound, Taylor-expand $\Psi$ about $0$ with the Lagrange remainder: there is
$u^\star\in[0,b^2+\mathcal V]$ with
$\Psi(b^2+\mathcal V)=\Psi(0)+\Psi'(0)(b^2+\mathcal V)+\tfrac12\Psi''(u^\star)(b^2+\mathcal V)^2
\le\kappa_U(b^2+\mathcal V)+\tfrac12 L_\Psi(b^2+\mathcal V)^2$, using $\Psi(0)=0$, $\Psi'(0)=\kappa_U$
and $\Psi''\le L_\Psi$. This is \eqref{eq:bv-bound}. \hfill$\square$ (part 1)

\par\smallskip\noindent\textbf{Step 3: the variance reduction factor $\theta$ is bias-free (the mechanism).}\quad 
We must show that replacing i.i.d.\ DP noise by anti-correlated noise (or adding an extra low-pass)
multiplies $\mathcal V$ by some $\theta(\lambda,T)$ \emph{without changing $b$}. The DP-MF mechanism
keeps the averaging weights $w_t$ and the privacy level fixed, so $\E[\widehat m_T]$ is unchanged
(the injected noise is zero-mean and the privacy-calibrated base std is rescaled only to preserve
sensitivity); hence $b$, and a fortiori the floor $\beta^2$, are invariant.

It remains to compute $\theta$. Inject $w_t^{\mathrm{noise}}=\xi_t-\lambda\xi_{t-1}$ with i.i.d.\
$\xi_t\sim\mathcal N(0,\nu^2)$. As a matrix mechanism, this is $n=Lz$ with $L$ lower-bidiagonal
($1$ on the diagonal, $-\lambda$ on the sub-diagonal); it realises the prefix-sum workload with
strategy $C_{\mathrm{strat}}=L^{-1}$, the lower-triangular Toeplitz matrix
$[1,\lambda,\lambda^2,\dots]$. The privacy-preserving per-step base std is inflated by the
\emph{max column $\ell_2$ norm} of $L^{-1}$,
\begin{equation}
\label{eq:kappa}
\kappa(\lambda,T)\;=\;\big\|L^{-1}\big\|_{\mathrm{col}}\;=\;\Big(\sum_{k=0}^{T-1}\lambda^{2k}\Big)^{1/2}\;=\;\sqrt{\frac{1-\lambda^{2T}}{1-\lambda^2}}\xrightarrow[T\to\infty]{}\frac{1}{\sqrt{1-\lambda^2}},
\end{equation}
the longest (first) column of $L^{-1}$. This is exactly \texttt{corr\_sensitivity(lam,steps)} in
\texttt{dp\_adaptive.py}; e.g.\ $\kappa(0.9,T{\to}\infty)=2.29$ (\emph{not} the naive $\sqrt{1+\lambda^2}=1.35$,
which is the column norm of $L$, under-noises, and would break privacy). We re-derived
\eqref{eq:kappa} independently (max-column-norm of $L^{-1}$) and confirmed it numerically against
$\|L^{-1}\|_{\mathrm{col}}$ over a grid of $(\lambda,T)$; it is $\sqrt{(1-\lambda^{2T})/(1-\lambda^2)}$,
\emph{not} $\sqrt{1+\lambda^2}$.

On the prefix-sum path ($w_t\equiv1/T$) the integrated injected noise telescopes:
$S_T:=\sum_{t=1}^{T}(\xi_t-\lambda\xi_{t-1})=\xi_T-\lambda\xi_0+(1-\lambda)\sum_{t=1}^{T-1}\xi_t$,
with i.i.d.\ $\xi_t$ of variance $\nu^2=\sigma^2\kappa^2$ at the matched privacy level. Hence
$\Var_{\mathrm{MF}}(S_T)=\sigma^2\kappa^2\big[\,1+\lambda^2+(1-\lambda)^2(T-1)\,\big]$, while the
i.i.d.\ control ($\lambda{=}0$, $\kappa{=}1$) gives $\Var_{\mathrm{iid}}(S_T)=\sigma^2 T$.
The exact finite-$T$ variance ratio is therefore
\begin{equation}
\label{eq:theta}
\theta(\lambda,T)\;=\;\frac{\Var_{\mathrm{MF}}(S_T)}{\Var_{\mathrm{iid}}(S_T)}\;=\;\frac{1-\lambda^{2T}}{1-\lambda^{2}}\cdot\frac{1+\lambda^2+(1-\lambda)^2(T-1)}{T}\;\xrightarrow[T\to\infty]{}\;\frac{1-\lambda}{1+\lambda}\ \in(0,1).
\end{equation}
We verified \eqref{eq:theta} both in closed form and by Monte-Carlo. \emph{Two caveats that the
asymptotic $(1-\lambda)/(1+\lambda)$ hides, and which matter here:} (a) at the experimental
$T\!\approx\!120$ the reduction is far milder than the asymptote---$\theta(0.9,120)\approx0.13$ (not
$0.05$), $\theta(0.95,120)\approx0.19$ (not $0.026$)---and is \emph{non-monotone} in $\lambda$
(minimised near $\lambda\!\approx\!0.9$, rising again toward $\lambda\!\to\!1$); (b) for small $T$
and large $\lambda$, $\theta(\lambda,T)$ can \emph{exceed} $1$ (e.g.\ $\theta(0.95,10)\approx1.27$),
i.e.\ the $\kappa$-inflation outweighs the telescoping cancellation and correlated noise
\emph{increases} the integrated variance. We restrict to the regime $\theta\in(0,1]$ (which holds
for the experimental $T\!\approx\!120$, all $\lambda\le0.99$); outside it the ``denoising'' premise
fails and part~2 is vacuous. In all cases $\theta(0,T)=1$ (i.i.d.\ control, reproduced
bit-for-bit), $\bar g$ is untouched, and $b$ is invariant. An extra causal low-pass and
noise-optimal momentum likewise multiply $\mathcal V$ by a factor $\le1$ while $b$ is fixed; they
are special cases of the same ``reduce $\mathcal V$, keep $b$'' operation. This is the precise sense
in which all three knobs act only on the variance.

\par\smallskip\noindent\textbf{Step 4: no-first-order-improvement corollary \eqref{eq:no-improve}--\eqref{eq:ratio}.}\quad 
Assume the signal-ceiling regime: $b^2\ge\beta^2>0$ with $b^2=\beta^2(1+o(1))$ ($\beta^2$
$\lambda$-free, Assumption~\ref{ass:floor}) and \eqref{eq:var-floor}, $\mathcal V\le\beta^2$. The
denoising replaces $\mathcal V$ by $\theta\mathcal V$ with $\theta\in(0,1]$ (Step 3), so
\[
\Delta U=U(\mathcal V)-U(\theta\mathcal V)=\Psi(b^2+\mathcal V)-\Psi(b^2+\theta\mathcal V).
\]
Since $\theta\le1$ and $\Psi$ is nondecreasing, $\Delta U\ge0$ (denoising never hurts in this
idealised model; the small \emph{observed} decreases at $\lambda>0$ are second-order
$\kappa$-inflation / window-mismatch / unamplified-accounting headwinds outside this monotone bound,
see Remark~\ref{rem:headwind}). By the mean value theorem there is
$u^\star\in[b^2+\theta\mathcal V,\ b^2+\mathcal V]$ with
\[
\Delta U=\Psi'(u^\star)\,(1-\theta)\,\mathcal V .
\]
By convexity ($\Psi''\le L_\Psi$), $\Psi'(u^\star)\le\Psi'(0)+L_\Psi u^\star=\kappa_U+L_\Psi u^\star$.
With $\mathcal V\le\beta^2$ and $b^2=\beta^2(1+o(1))$ we have
$u^\star\le b^2+\mathcal V\le 2\beta^2(1+o(1))$, hence $u^\star\mathcal V\le2\beta^4(1+o(1))$. Thus
\[
\Delta U\le(\kappa_U+L_\Psi u^\star)(1-\theta)\mathcal V\le\kappa_U(1-\theta)\mathcal V+L_\Psi u^\star\mathcal V
\;\le\;\kappa_U(1-\theta)\mathcal V+2L_\Psi\beta^4\,(1+o(1)),
\]
which is \eqref{eq:no-improve}. (For the \emph{derived} $\Psi$ of Lemma~\ref{lem:util}, $L_\Psi=0$
and this is the exact equality $\Delta U=\tfrac\eta2(1-\theta)\mathcal V$, with no $O(\beta^4)$
term.) The irreducible loss on which $\Delta U$ sits is $\Psi(b^2)\ge\Psi(\beta^2)\ge\kappa_U\beta^2$
(the last by convexity and $\Psi(0)=0$: $\Psi(\beta^2)\ge\Psi(0)+\Psi'(0)\beta^2=\kappa_U\beta^2$).
Dividing,
\[
\frac{\Delta U}{\Psi(b^2)}\;\le\;\frac{\kappa_U(1-\theta)\mathcal V+2L_\Psi\beta^4}{\kappa_U\beta^2}\;=\;(1-\theta)\frac{\mathcal V}{\beta^2}\;+\;\frac{2L_\Psi}{\kappa_U}\beta^2 .
\]
Using $\mathcal V\le\beta^2$ gives $(1-\theta)\frac{\mathcal V}{\beta^2}\le(1-\theta)$, and the
second term is $O(\beta^2)\to0$ as the floor is approached (and is identically $0$ for the
descent-lemma $\Psi$), proving \eqref{eq:ratio}. In the strict signal-ceiling limit where averaging
has driven the variance far below the floor, $\mathcal V/\beta^2\to0$, the bound gives
$\Delta U/\Psi(b^2)\to0$ \emph{for every} $\theta\in(0,1]$: even $\theta\to0$ (variance annihilated)
cannot move the relative utility, because the loss is pinned at the bias floor $\Psi(\beta^2)$. This
is the no-improvement corollary. \hfill$\square$ (part 2)

\par\smallskip\noindent\textbf{Step 5: non-transfer of DP-MF (part 3).}\quad 
Now take the opposite regime $v_{\mathrm{true}}\gg\Phi$, $\rho\to0$ (large true gradients, e.g.\
from-scratch DP-SGD where the signal is well above the per-step noise floor). Here the clipping bias
is the only bias and the below-floor deficit vanishes, so $\beta^2$ is negligible and, because the
per-step noise is large relative to the signal averaged over few effective steps,
$\mathcal V\gg b^2\approx\beta^2\approx0$, i.e.\ the estimator is \emph{variance-limited}. Then
$u^\star=b^2+\Theta(\mathcal V)=\Theta(\mathcal V)$, and provided $\Psi$ is locally affine on the
relevant scale (or $L_\Psi\mathcal V\ll\kappa_U$), the mean-value form gives
\[
\Delta U=\Psi'(u^\star)(1-\theta)\mathcal V=\kappa_U(1-\theta)\mathcal V\,(1+o(1)),
\]
a \emph{genuine first-order} gain proportional to the variance-reduction fraction $(1-\theta)$,
which for correlated noise is $1-\theta\to\tfrac{2\lambda}{1+\lambda}$ (large $T$) and can approach
$1$. Relative to the (now small) irreducible loss this is an unbounded improvement factor:
DP-MF strictly helps. Comparing Steps~4 and 5: the \emph{same} mechanism (multiply $\mathcal V$ by
$\theta$, keep $b$) produces a first-order utility gain iff $\mathcal V\gg\beta^2$ (variance-limited,
$\rho\to0$) and is first-order inert iff $\mathcal V\lesssim\beta^2$ (signal/bias-limited,
$\rho\to1$). Since gentle LLM DP-LoRA fine-tuning is measured at $\rho=1.00$--$1.07$ (signal below
the per-step floor, $v_{\mathrm{true}}\approx0$) while from-scratch DP training operates at
$\rho\ll1$, the DP-MF / correlated-noise improvement that is provable in the latter does not
transfer to the former. This is the formal content of the non-transfer claim. \hfill$\square$ (part 3)

\par\smallskip\noindent\textbf{Step 6: consistency with the empirical record.}\quad 
The theorem predicts $\Delta U/\Psi(b^2)\le(1-\theta)+O(\beta^2)$ with the leading term vanishing
once $\mathcal V\ll\beta^2$, i.e.\ near-ties with possible small sign-indefinite wiggles at the bar.
This matches every measured outcome at $\rho\approx1$: DP-CorrMom on Adam-momentum
($\lambda{=}0$: $55.67$ vs $\lambda{=}0.9/0.95$: $54.35/54.10$), $\beta_1{=}0$
($\lambda\in\{0,0.9,0.95,0.99\}=56.01/54.35/54.26/52.19$, monotone non-improving), the matched
plain-SGD workload (\texttt{dp-corrsgd} lr$10$: $\lambda{=}0$: $55.70$ vs $\lambda{=}0.95$: $55.35$),
and the first-moment low-pass on RoBERTa/MNLI ($0.35$ chance vs $0.71$--$0.75$ plain DP-Adam): in no
case does first-moment denoising beat the $\lambda{=}0$ / tuned, amplified DP-Adam reference
($56.22$ at $\varepsilon{=}6.404$, less budget). The theorem explains \emph{why}: at $\rho\approx1$
the loss is pinned at $\Psi(\beta^2)$ and the variance term is already below the floor, so the
$(1-\theta)\mathcal V/\beta^2$ first-order lever is empty. \hfill$\blacksquare$

\begin{remark}[Why the empirical sign can be negative, not just zero]\label{rem:headwind}
The idealised model gives $\Delta U\ge0$ (denoising weakly helps). The small observed
\emph{decreases} at $\lambda>0$ are second-order headwinds outside the bias--variance abstraction:
(i) the privacy-preserving $\kappa(\lambda,T)$-inflation of the per-step base std
(\eqref{eq:kappa}; $2.29\times$ at $\lambda{=}0.9$) raises the \emph{per-step} noise in the
$\sqrt{\hat v}$ denominator faster than the prefix-sum cancellation lowers the integrated noise when
the averaging window is short---Adam's $\approx$10-step EMA does \emph{not} match the full prefix
sum the strategy is matched to, and indeed $\theta(\lambda,T)$ is far from its $(1-\lambda)/(1+\lambda)$
asymptote at the realised $T$ (Step~3); and (ii) the unamplified (route-A) accountant that
correlated noise forces spends more budget than the Poisson-amplified DP-Adam at equal steps. Both
push $\Delta U$ slightly negative once the first-order $(1-\theta)\mathcal V/\beta^2$ term is
$\approx0$, exactly as the signal-ceiling predicts a vanishing first-order benefit on which these
second-order costs then dominate.
\end{remark}

\begin{remark}[Tightness and scope]\label{rem:scope}
The bound is tight in the limit: at $\mathcal V=\beta^2$, $\theta\to0$ and $L_\Psi=0$ (the
descent-lemma $\Psi$), \eqref{eq:ratio} gives $\Delta U/\Psi(b^2)\le1$, i.e.\ denoising could in
principle halve the on-floor loss when the variance equals the bias floor; the no-improvement
statement is therefore specifically about the \emph{strict} ceiling $\mathcal V\ll\beta^2$ that
basic momentum already achieves ($\sqrt T$ variance reduction over $T\gg1$ steps drives $\mathcal V$
well below $\beta^2$). The scope is LoRA-scale, $\rho\approx1$ fine-tuning; the dichotomy in part 3
delineates exactly when the conclusion flips.
\end{remark}